%% ── Section 1: Introduction ─────────────────────────────────
\section{Introduction}

Diabetic Retinopathy (DR) is a common complication of diabetes and a leading
cause of preventable blindness in working-age adults. It damages the small blood
vessels of the retina, which appears in fundus photographs as microaneurysms,
haemorrhages, hard exudates and, in advanced cases, the growth of new vessels.
In countries such as Bangladesh, where more than 13.1~million people live with
diabetes~\cite{idf2021}, the burden is made worse by a shortage of retinal
specialists outside the major cities. DR is graded on a five-level ordinal scale:
Grade~0 (no DR), Grade~1 (mild non-proliferative DR, NPDR), Grade~2 (moderate
NPDR), Grade~3 (severe NPDR), and Grade~4 (proliferative DR). Referable disease
(Grade~$\geq$~2) has to be detected early to prevent vision loss, which makes
automated screening attractive.

When all images are pooled in one place, deep convolutional networks grade DR
accurately. On the imbalanced DDR dataset~\cite{li2019ddr}, an EfficientNet-B3
model reaches a Quadratic Weighted Kappa (QWK) of 0.8494.
% TODO(Day 4): new centralised QWK (results/centralised_test.json)
QWK is the standard metric for ordinal grading because it penalises an error by
the square of its distance: predicting Grade~0 for a Grade~4 image costs far more
than predicting Grade~3. Pooling patient images from several hospitals, however,
is restricted by privacy regulations such as HIPAA and GDPR. Federated Learning
(FL)~\cite{mcmahan2017fedavg} avoids this problem: the hospitals train a shared
model by exchanging model weights, and the images never leave each site.

The main difficulty in clinical FL is that hospitals see different patients, so
their data are not identically distributed (non-IID). Screening clinics see
mostly healthy (Grade~0) eyes, referral centres see many severe cases
(Grades~3--4), and rural clinics see cases in between. Under extreme Dirichlet
skew ($\alpha=0.1$), standard FedAvg is observed to undergo catastrophic
forgetting, crashing from a peak QWK of 0.7993 down to 0.4064 in a single
communication round. FedProx~\cite{li2020fedprox}, which limits how far each local
model can move away from the global model, reduces but does not prevent this
drop (down to 0.5963).
% TODO(Day 4): replace 0.7993 / 0.4064 / 0.5963 with the new multi-seed results,
% and soften "in a single communication round" if the new runs do not show it.
Both methods weight each hospital by its number of images, so a large clinic that
sees mostly one grade can dominate the global model.

In addition, the standard cross-entropy loss ignores the order of the grades:
mistaking Grade~4 for Grade~0 costs the same as mistaking it for Grade~3. Ordinal
losses~\cite{toledocortes2022dqor}--[\citen{polat2024cdwce}] address this in
centralised training, but their effect on client drift in FL has received little
attention.

To address both problems, \textbf{OrdinalFed} is proposed: a federated framework
for ordinal medical grading with two components.
\begin{enumerate}[leftmargin=*]
  \item \textbf{Distance-aware OrdinalLoss:} a local training loss that adds a
  squared expected-grade penalty to class-weighted cross-entropy, so that
  predictions far from the true grade are penalised more heavily.
  \item \textbf{QWK-weighted aggregation with an adaptive $\beta$:} the server
  scores each client model on a small balanced set $\mathcal{V}$ of 50 images and
  gradually shifts the aggregation weights from image counts towards these QWK
  scores, which limits the influence of clients whose updates degrade on
  $\mathcal{V}$.
\end{enumerate}
